\section{Results}

\subsection{Cohort Characteristics and Data Quality}

We analyzed 730 participants from the Parkinson's Progression Markers Initiative (PPMI) comprising 536 individuals with Parkinson's disease (PD) and 194 prodromal participants exhibiting risk markers for future PD diagnosis. The PD cohort had a mean age of 61.3 $\pm$ 9.7 years, with 37.1\% female participants, while the prodromal cohort was slightly older (63.8 $\pm$ 8.2 years) with 35.1\% female participants (Table 1). At baseline, the PD cohort demonstrated a mean MDS-UPDRS Part III motor score of 21.4 $\pm$ 10.2 points and Montreal Cognitive Assessment (MoCA) score of 27.1 $\pm$ 2.4, indicating mild motor symptoms with preserved cognition. The prodromal cohort showed minimal motor symptoms (UPDRS-III: 5.2 $\pm$ 3.8) and normal cognition (MoCA: 28.3 $\pm$ 1.9).

\paragraph{Imaging and Genetic Biomarkers}
Dopamine transporter imaging (DaTscan SPECT) revealed reduced striatal binding ratios (SBR) in the PD cohort (putamen SBR: 1.52 $\pm$ 0.38) compared to prodromal participants (1.89 $\pm$ 0.35, $p < 0.001$), consistent with dopaminergic denervation. Structural MRI data from FreeSurfer analysis provided regional cortical thickness measurements across 68 brain regions (34 per hemisphere). Genetic screening identified 7.8\% of PD participants and 9.3\% of prodromal participants carrying pathogenic \textit{GBA} mutations, 4.3\% and 5.7\% carrying \textit{LRRK2} mutations, and 31.2\% and 28.9\% carrying at least one \textit{APOE} $\varepsilon$4 allele, respectively.

\paragraph{Data Quality and Preprocessing}
Initial data quality assessment revealed missing values ranging from 2.1\% (demographic variables) to 18.3\% (CSF biomarkers). We applied K-nearest neighbors (KNN) imputation with k=5 for missing values below 20\% and excluded features with >20\% missingness. After quality control, we retained 87 features across four modalities: clinical assessments (32 features), imaging biomarkers (42 features from DaTscan and sMRI), genetic variants (8 features), and CSF markers (5 features). Data scaling was performed using standardization (z-score normalization) separately for each modality to account for different measurement scales.

\paragraph{Patient Similarity Graph Characteristics}
The constructed patient similarity graphs demonstrated clinically meaningful structure. For the PD cohort, the k-nearest neighbor graph (k=10) showed a mean node degree of 10.0 $\pm$ 0.0 by construction, with an average clustering coefficient of 0.32 $\pm$ 0.14, indicating moderate local connectivity. Graph homophily, measuring the tendency of connected nodes to share the same diagnostic label, was 0.73 $\pm$ 0.08, significantly above random expectation ($p < 0.001$). The prodromal cohort graph exhibited similar topological properties (clustering coefficient: 0.29 $\pm$ 0.12, homophily: 0.68 $\pm$ 0.09).

Edge weight analysis revealed that clinical similarity contributed most strongly to patient connections (mean contribution: 38.2\%), followed by imaging features (31.7\%), genetic similarity (18.4\%), and CSF biomarkers (11.7\%). This distribution reflects both the higher dimensionality of clinical assessments and their greater variability across the cohort. Graph connectivity analysis showed that 98.7\% of PD participants and 97.4\% of prodromal participants were part of the largest connected component, ensuring that graph-based learning could effectively propagate information across the cohort.

\paragraph{Longitudinal Data Characteristics}
For the Phase 4 longitudinal subtyping analysis, we utilized repeated measurements from PD participants with at least 3 visits over a minimum 2-year follow-up period (n=536, mean follow-up: 4.2 $\pm$ 1.8 years, range: 2.0-10.0 years). The dataset included a total of 2,847 clinical visits (mean: 5.3 visits per participant). For Phase 5 prodromal progression prediction, the prodromal cohort (n=194) was followed for a mean of 3.6 $\pm$ 1.5 years (range: 1.0-8.0 years), during which 47 participants (24.2\%) converted to clinical PD diagnosis, 89 participants (45.9\%) remained prodromal, and 58 participants (29.9\%) were censored.
